\section{Introduction}
\label{sec:introduction}

Storage maintenance is increasingly driven by signals that arrive before a hard failure: SMART counters, I/O errors, latency changes, imbalance, corruption detected by scrubbing, and workload pressure. These signals create an opportunity to move from failure handling after the fact to continuous maintenance. They also create a control problem. A transient warning should not automatically evict a healthy device, an aggressive scrub should not overwhelm foreground I/O, and a controller must not perform a destructive transition when redundancy is already insufficient.

Prior work has studied automatic storage configuration, disk-failure prediction, and proactive repair as important individual mechanisms \citep{anderson2002hippodrome,lu2020smarter,li2022fastpr}. We focus on the layer between prediction and repair: how a file system can safely turn imperfect observations into maintenance actions while preserving availability and bounding interference with foreground work.

We study this problem through \autopilot, the autonomous maintenance controller in \argosfs. The design treats maintenance as a persistent feedback loop rather than a collection of independent timers. Each iteration gathers observations, updates durable risk memory, evaluates safety and policy constraints, selects a bounded action, records the outcome, and uses recent foreground pressure and action utility to constrain subsequent work.

The paper is structured around three contributions:
\begin{itemize}[leftmargin=*]
  \item We formulate \emph{safe closed-loop maintenance} as a control problem under uncertain health observations, explicit redundancy constraints, and shared foreground/background I/O resources.
  \item We design and implement \autopilot, which unifies observe, scrub, rebalance, and drain decisions with persistent evidence, confirmation and cooldown gates, conflict-aware mutation stopping, and foreground-aware budgets.
  \item We define a reproducible evaluation that separates correctness from effectiveness: large trace-driven decision tests, real block-fault injection, mounted file-system workloads, and ablations against reactive, periodic, fixed-threshold, and oracle policies.
\end{itemize}

\Cref{sec:background} states the problem and metrics. \Cref{sec:design} presents the controller, \cref{sec:implementation} maps the design to \argosfs, and \cref{sec:evaluation} defines the experiments used to validate the claims.
